% Unit 108 terjemahan Bahasa Indonesia dari chapter8.tex baris 282--453.
% Sumber: Methods of Algebra, Volume 2, commit
% 9a5803ff77dd3257484cb177f851a73770a59dd3 (CC BY 4.0).
% stable-unit-id: o014.aljabr2.chapter8.nerves-of-categories
% Cakupan: Bagian 8.3 lengkap, ``Contoh: Saraf Kategori'';
% berhenti tepat sebelum Bagian 8.4 pada baris 454.

% segment-id: o014.li2.chapter8.nerves-of-categories.h001
\section{Contoh: Saraf Kategori}\label{sec:nerves}
\hypertarget{o014.aljabr2.chapter8.nerves-of-categories}{ }
% segment-id: o014.li2.chapter8.nerves-of-categories.p001
Bagian ini memperkenalkan saraf kategori dan, melaluinya, ruang klasifikasi
grup. Keduanya merupakan contoh penting himpunan simplisial.

% segment-id: o014.li2.chapter8.nerves-of-categories.eg001
\begin{example}[Saraf kategori]\label{eg:nerve-cat}
	\index{saraf@saraf (nerve)}
	\index[sym1]{NC@$\mathrm{N}\mathcal{C}$}
	Misalkan $\mathcal{C}$ kategori kecil. Definisikan funktor
	% segment-id: o014.li2.chapter8.nerves-of-categories.q001
	\[ \simpDelta^{\opp} \to \cate{Set}, \quad [n] \mapsto \left\{ \text{semua funktor } [n] \to \mathcal{C} \right\}. \]
	Sebagai himpunan simplisial, funktor ini ditulis
	$\mathrm{N}\mathcal{C}$ dan disebut \emph{saraf} $\mathcal{C}$. Memberikan
	elemen $\mathrm{N}\mathcal{C}_n$ sama artinya dengan memberikan funktor
	$[n]\to\mathcal{C}$, yakni rantai morfisme dalam $\mathcal{C}$
	% segment-id: o014.li2.chapter8.nerves-of-categories.q002
	\[ (f_1, \ldots, f_n): C_0 \xrightarrow{f_1} C_1 \xrightarrow{f_2} \cdots \xrightarrow{f_n} C_n. \]
	Jadi, $\mathrm{N}\mathcal{C}_0$ dapat diidentifikasi dengan
	$\Obj(\mathcal{C})$, dan $\mathrm{N}\mathcal{C}_1$ dengan
	$\Mor(\mathcal{C})$.

	Morfisme muka $d_i:\mathrm{N}\mathcal{C}_n\to
	\mathrm{N}\mathcal{C}_{n-1}$ dan morfisme degenerasi
	$s_j:\mathrm{N}\mathcal{C}_n\to\mathrm{N}\mathcal{C}_{n+1}$ diberikan
	oleh
	% segment-id: o014.li2.chapter8.nerves-of-categories.q003
	\begin{align*}
		d_i(f_1, \ldots, f_n) & = \begin{cases}
			(f_2, \cdots, f_n), & i = 0 \\
			(\ldots, f_{i+1} f_i, \ldots), & 0 < i < n, \\
			(f_1, \ldots, f_{n-1}), & i = n,
		\end{cases} \\
		s_j(f_1, \ldots, f_n) & = \begin{cases}
			\left( \identity_{C_0}, f_1, \ldots, f_n \right), & j = 0 \\
			(\ldots, f_j, \identity_{C_j}, f_{j+1}, \ldots ), & 0 < j < n \\
			(f_1, \ldots, f_n, \identity_{C_n}), & j = n.
		\end{cases}
	\end{align*}
\end{example}

% segment-id: o014.li2.chapter8.nerves-of-categories.r001
\begin{remark}
	Dengan menguraikan definisi, terlihat bahwa himpunan simplisial
	$\mathrm{N}(\mathcal{C}^{\opp})$ dan $\mathrm{N}\mathcal{C}$ dihubungkan
	oleh dualitas pembalikan urutan dalam
	Catatan~\ref{rem:simpDelta-order-reversal}. Argumen terperinci diserahkan
	kepada pembaca.
\end{remark}

% segment-id: o014.li2.chapter8.nerves-of-categories.p002
Saraf merealisasikan kategori secara konkret melalui data kombinatorial/topologis
dan memuat seluruh informasi kategori asal. Hal ini merupakan isi
Proposisi~\ref{prop:nerve-ff}, yang akan segera dibuktikan. Mula-mula kita
akan mencirikan himpunan simplisial mana yang merupakan saraf; argumennya menarik.

% segment-id: o014.li2.chapter8.nerves-of-categories.pr001
\begin{proposition}[Pencirian saraf]\label{prop:inner-horn-filling}
	Misalkan $X$ himpunan simplisial. Pernyataan berikut ekuivalen.
	\begin{enumerate}[(i)]
		\item Terdapat kategori kecil $\mathcal{C}$ sedemikian sehingga
		$\mathrm{N}\mathcal{C}\simeq X$.
		\item $X$ mempunyai sifat pengisian tanduk dalam yang unik: untuk setiap
		bilangan bulat $0<i<n$ dan morfisme
		$\sigma':\Lambda^n_i\to X$ dalam $\cate{sSet}$ (objek
		$\Lambda^n_i$ seperti ini disebut tanduk dalam), terdapat tepat satu
		$\sigma:\Delta^n\to X$ yang memperluas $\sigma'$.
	\end{enumerate}
\end{proposition}
% segment-id: o014.li2.chapter8.nerves-of-categories.pf001
\begin{proof}
	Pertama, kita buktikan (i) $\implies$ (ii). Misalkan
	$X=\mathrm{N}\mathcal{C}$ dan $0<i<n$; kita ingin memperluas
	$\sigma':\Lambda^n_i\to X$ ke $\Delta^n$. Untuk setiap
	$0\leq k\leq n$ (atau $0<k\leq n$), morfisme
	$\Delta^0\xrightarrow{\{k\}}\Delta^n$ (atau
	$\Delta^1\xrightarrow{\{k-1,k\}}\Delta^n$) terfaktorkan melalui
	$\Lambda^n_i$. Tuliskan citranya di bawah $\sigma'$ sebagai
	$C_k\in X_0=\Obj(\mathcal{C})$ (atau
	$[g_k:C_{k-1}\to C_k]\in X_1=\Mor(\mathcal{C})$). Dengan demikian,
	diperoleh elemen $X_n$
	% segment-id: o014.li2.chapter8.nerves-of-categories.q004
	\[ C_0 \xrightarrow{g_1} C_1 \xrightarrow{g_2} \cdots \xrightarrow{g_n} C_n. \]
	Tuliskan morfisme $\Delta^n\to X$ yang bersesuaian sebagai $\sigma$.
	Konstruksi ini sekaligus menunjukkan bahwa $\sigma$ merupakan satu-satunya
	kemungkinan perluasan $\sigma'$.
	Karena
	$\Lambda^n_i=\bigcup_{j\neq i}\Image[\mathrm{d}^j:
	\Delta^{n-1}\to\Delta^n]$, untuk membuktikan pernyataan (ii) cukup diperiksa bahwa
	% segment-id: o014.li2.chapter8.nerves-of-categories.q005
	\begin{equation}\label{eqn:inner-horn-filling-aux0}
		\sigma \circ \mathrm{d}^j = \sigma' \circ \mathrm{d}^j : \Delta^{n-1} \to X, \quad j \neq i .
	\end{equation}

	Menurut definisi saraf, untuk membuktikan
	\eqref{eqn:inner-horn-filling-aux0} cukup memeriksa hal berikut. Untuk
	setiap $j\neq i$ dan setiap pasangan elemen
	bertetangga $h<k$ dalam deret $0,\ldots,\widehat{j},\ldots,n$ (simbol
	$\widehat{j}$ berarti bahwa $j$ dihapus), penarikan balik $\sigma$ dan
	$\sigma'$ sepanjang
	$\Delta^1\xrightarrow{\{h,k\}}\Lambda^n_i\subset\Delta^n$ harus sama.
	\begin{itemize}
		\item Jika $k=h+1$, kesamaan ini dijamin oleh konstruksi $\sigma$.
		Khususnya, \eqref{eqn:inner-horn-filling-aux0} selalu berlaku bila
		$j\in\{0,n\}$.
		\item Misalkan $(h,k)=(j-1,j+1)$. Jika $n=2$, keadaan ini tidak mungkin.
		Jika $n>2$, berlaku $j-1>0$ atau $j+1<n$. Bila $j-1>0$,
		$\{j-1,j+1\}$ terfaktorkan sebagai
		% segment-id: o014.li2.chapter8.nerves-of-categories.q006
		\[ \Delta^1 \xrightarrow{\{j-2, j\}} \Delta^{n-1} \xrightarrow{\mathrm{d}^0 = \{1, \ldots, n\}} \Lambda^n_i \subset \Delta^n. \]
		Menurut kasus $j=0$ dari
		\eqref{eqn:inner-horn-filling-aux0} yang sudah diketahui, kedua penarikan
		balik itu sama. Serupa dengan itu, bila $j+1<n$, persoalan
		tereduksi pada kasus $j=n$ yang telah diketahui.
	\end{itemize}

	Sekarang kita buktikan (ii) $\implies$ (i). Definisikan kategori $\mathcal{C}$
	dengan $\Obj(\mathcal{C}):=X_0$ dan, untuk sebarang $C,C'\in X_0$,
	% segment-id: o014.li2.chapter8.nerves-of-categories.q007
	\[ \Hom_{\mathcal{C}}(C, C') := \left\{ f \in X_1: d_1(f) = C, \; d_0(f) = C' \right\}. \]
	Gunakan pemetaan degenerasi $s_0:X_0\to X_1$ untuk mendefinisikan morfisme
	identitas $\identity_C$ sebagai $s_0(C)$. Selanjutnya, kita juga
	menggambarkan morfisme $f$ sebagai $0\xrightarrow{f}1$ untuk menekankan
	bahwa ia bersesuaian dengan simpleks-$1$ $[1]=\{0,1\}\to X$.

	Jika $d_0(f)=d_1(g)$, definisikan komposisi sebagai
	$gf:=d_1(\sigma)$, dengan $\sigma:\Delta^2\to X$ perluasan unik dari
	% segment-id: o014.li2.chapter8.nerves-of-categories.q008
	\[\begin{tikzcd}[column sep=small]
		& 1 \arrow[rd, "g"] & \\
		0 \arrow[ru, "f"] & & 2
	\end{tikzcd}: \Lambda^2_1 \to X
	\quad \text{yakni}\quad d^0(\sigma) = g, \quad d^2(\sigma) = f. \]

	Sifat $f\circ\identity_C=f$ dan $\identity_{C'}\circ f=f$ masing-masing
	dibuktikan oleh elemen-elemen $X_2$ berikut.
	% segment-id: o014.li2.chapter8.nerves-of-categories.q009
	\[ \begin{tikzcd}[row sep=large]
		& 1 \arrow[rd, bend left, "{\identity_{C'} = d_0 s_1(f)}"] & \\
		0 \arrow[rr, "{}" name=A, "{f = d_1 s_1(f)}"'] \arrow[ru, bend left, "{f = d_2 s_1(f)}"] & & 2 \arrow[phantom, from=1-2, to=A, "{s_1(f)}" description]
	\end{tikzcd} \quad \begin{tikzcd}[row sep=large]
		& 1 \arrow[rd, bend left, "{f = d_0 s_0(f)}"] & \\
		0 \arrow[rr, "{}" name=A, "{f = d_1 s_0(f)}"'] \arrow[ru, bend left, "{\identity_C = d_2 s_0(f)}"] & & 2 \arrow[phantom, from=1-2, to=A, "{s_0(f)}" description]
	\end{tikzcd}\]

	Untuk asosiativitas $h(gf)=(hg)f$, konstruksikan
	$\sigma':\Lambda^3_2\to X$ dengan ketiga mukanya masing-masing
	% segment-id: o014.li2.chapter8.nerves-of-categories.q010
	\[\begin{tikzcd}[column sep=small]
		& 1 \arrow[rd, "g"] & \\
		0 \arrow[ru, "f"] \arrow[rr, "gf"'] & & 2
	\end{tikzcd} \quad \begin{tikzcd}[column sep=small]
		& 2 \arrow[rd, "h"] & \\
		1 \arrow[ru, "g"] \arrow[rr, "hg"'] & & 3
	\end{tikzcd} \quad \begin{tikzcd}[column sep=small]
		& 2 \arrow[rd, "h"] & \\
		0 \arrow[ru, "gf"] \arrow[rr, "h(gf)"'] & & 3.
	\end{tikzcd}\]
	Tanduk ini mempunyai perluasan unik $\sigma:\Delta^3\to X$ dengan
	$d^2(\sigma)\in X_2$ berbentuk
	% segment-id: o014.li2.chapter8.nerves-of-categories.q011
	\[\begin{tikzcd}[column sep=small]
		& 1 \arrow[rd, "hg"] & \\
		0 \arrow[ru, "f"] \arrow[rr, "h(gf)"'] & & 3.
	\end{tikzcd} \]
	Namun, menurut konstruksi, simpleks-$2$ ini juga menentukan komposisi $f$
	dan $hg$, sehingga membuktikan $(hg)f=h(gf)$.

	Jadi, $\mathcal{C}$ adalah kategori. Terdapat morfisme kanonik
	$X\to\mathrm{N}\mathcal{C}$: dari suatu $\Delta^n\to X$, tarik balik
	sepanjang setiap $\Delta^1\xrightarrow{\{j,j+1\}}\Delta^n$ untuk
	memperoleh rantai morfisme. Dari konstruksinya, jelas bahwa
	$X_n\to\mathrm{N}\mathcal{C}_n$ bijektif untuk $n=0,1$. Bila $n\geq2$,
	pilih $0<i<n$ dan tinjau diagram komutatif
	% segment-id: o014.li2.chapter8.nerves-of-categories.q012
	\[\begin{tikzcd}
		\Hom(\Delta^n, X) \arrow[d] \arrow[r] & \Hom(\Delta^n, \mathrm{N}\mathcal{C}) \arrow[d] \\
		\Hom(\Lambda^n_i, X) \arrow[r] & \Hom(\Lambda^n_i, \mathrm{N}\mathcal{C}).
	\end{tikzcd}\]
	Menurut asumsi dan sifat pengisian tanduk dalam yang unik, kedua anak panah
	vertikal merupakan bijeksi. Selain itu, $\Lambda^n_i$ dapat dinyatakan sebagai
	$\varinjlim$ suatu keluarga $\Delta^{n-1}$ dan $\Delta^{n-2}$ (yakni
	penempelan $n$ muka), sehingga melalui induksi anak panah pada baris kedua
	juga merupakan bijeksi. Selesailah bukti.
\end{proof}

% segment-id: o014.li2.chapter8.nerves-of-categories.p003
Perhatikan pula bahwa dalam argumen (ii) $\implies$ (i), jika terdapat
kategori kecil $\mathcal{C}_0$ dengan $X=\mathrm{N}(\mathcal{C}_0)$, maka
kategori $\mathcal{C}$ yang dikonstruksi dalam bukti tepat sama dengan
$\mathcal{C}_0$.

% segment-id: o014.li2.chapter8.nerves-of-categories.p004
Tuliskan $\cate{Cat}$ untuk kategori semua kategori kecil, dengan funktor
sebagai morfisme. Setiap funktor $F:\mathcal{C}\to\mathcal{C}'$ menginduksi
morfisme $\mathrm{N}F:\mathrm{N}\mathcal{C}\to
\mathrm{N}\mathcal{C}'$ dalam $\cate{sSet}$ yang memetakan
$(f_1,\ldots,f_n)$ ke $(Ff_1,\ldots,Ff_n)$. Dengan demikian diperoleh funktor
saraf $\mathrm{N}:\cate{Cat}\to\cate{sSet}$.

% segment-id: o014.li2.chapter8.nerves-of-categories.pr002
\begin{proposition}\label{prop:nerve-ff}
	Funktor saraf $\mathrm{N}:\cate{Cat}\to\cate{sSet}$ penuh setia.
\end{proposition}
% segment-id: o014.li2.chapter8.nerves-of-categories.pf002
\begin{proof}
	Bukti Proposisi~\ref{prop:inner-horn-filling} telah menunjukkan cara
	merekonstruksi morfisme dan komposisinya dari saraf kategori. Karena itu,
	morfisme $\mathrm{N}(\mathcal{C})\to\mathrm{N}(\mathcal{C}')$ secara alami
	menginduksi funktor $\mathcal{C}\to\mathcal{C}'$. Mudah dilihat bahwa
	konstruksi ini saling invers dengan konstruksi yang diinduksi oleh $\mathrm{N}$.
\end{proof}

% segment-id: o014.li2.chapter8.nerves-of-categories.eg002
\begin{example}[Ruang klasifikasi]\label{eg:classifying-space}
	\index{ruang klasifikasi@ruang klasifikasi (classifying space)}
	\index[sym1]{BGamma@$\mathrm{B}\Gamma$}
	\index[sym1]{EGamma@$\mathrm{E}\Gamma$}
	Misalkan $\Gamma$ monoid. Definisikan kategori $\mathcal{B}\Gamma$ dan
	$\mathcal{E}\Gamma$ sebagai berikut:
	% segment-id: o014.li2.chapter8.nerves-of-categories.q013
	\begin{center}\begin{tabular}{|c|c|c|c|}\hline
		Kategori & Himpunan objek & Morfisme & Komposisi morfisme \\ \hline
		$\mathcal{B}\Gamma$ & $\{\star\}$ & $\End_{\mathcal{B}\Gamma}(\star)=\Gamma$ & perkalian dalam $\Gamma$ \\ \hline
		$\mathcal{E}\Gamma$ & $\Gamma$ & $\Hom_{\mathcal{E}\Gamma}(g,g')=\{h\in\Gamma:hg=g'\}$ & perkalian dalam $\Gamma$ \\ \hline
	\end{tabular}\end{center}

	Definisikan saraf
	$\mathrm{B}\Gamma:=\mathrm{N}((\mathcal{B}\Gamma)^{\opp})$ dan
	$\mathrm{E}\Gamma:=\mathrm{N}((\mathcal{E}\Gamma)^{\opp})$.
	Himpunan $\mathrm{B}\Gamma$ disebut ruang klasifikasi $\Gamma$. Mengambil
	Penerapan $(\cdots)^{\opp}$ pada dasarnya membalik anak panah dalam definisi saraf
	tanpa mengubah labelnya. Karena itu,
	% segment-id: o014.li2.chapter8.nerves-of-categories.q014
	\begin{equation*}
		\begin{array}{|c|c|c|} \hline
		\text{Himpunan} & \text{Elemen} & \text{Rantai morfisme terkait dalam }\mathcal{B}\Gamma\text{ atau }\mathcal{E}\Gamma \\ \hline
		(\mathrm{B}\Gamma)_n = \Gamma^n & (g_1, \ldots, g_n) & \star \xleftarrow{g_1} \cdots \leftarrow \star \xleftarrow{g_n} \star \\ \hline
		(\mathrm{E}\Gamma)_n = \Gamma^{n+1} & (g_1, \ldots, g_{n+1}) & g_1 \cdots g_{n+1} \xleftarrow{g_1} \cdots \leftarrow g_n g_{n+1} \xleftarrow{g_n} g_{n+1} \\ \hline
		\end{array}
	\end{equation*}
	Dengan langsung memasukkan definisi, untuk setiap $n\geq0$ diperoleh
	% segment-id: o014.li2.chapter8.nerves-of-categories.q015
	\begin{equation*}
		\mathrm{B}\Gamma: \qquad \begin{aligned}
			d_i(g_1, \ldots, g_n) & =
			\begin{cases}
				(g_2, \ldots, g_n), & i = 0 \\
				(\ldots, g_i g_{i+1} , \ldots), & 0 < i < n \\
				(g_1, \ldots, g_{n-1}), & i = n,
			\end{cases} \\
			s_j(g_1, \ldots, g_n) & =
			\begin{cases}
				(1, g_1, \ldots, g_n), & j = 0 \\
				(\ldots, g_j, 1, g_{j+1}, \ldots), & 0 < j < n \\
				(g_1, \ldots, g_n, 1), & j = n,
			\end{cases}
		\end{aligned}
	\end{equation*}
	dan
	% segment-id: o014.li2.chapter8.nerves-of-categories.q016
	\begin{equation*}
		\mathrm{E}\Gamma: \qquad \begin{aligned}
			d_i(g_1, \ldots, g_{n+1}) & =
			\begin{cases}
				(g_2, \ldots, g_{n+1}), & i = 0 \\
				(g_1, \ldots, g_i g_{i+1} , \ldots), & 0 < i \leq n,
			\end{cases} \\
			s_j(g_1, \ldots, g_{n+1}) & =
			\begin{cases}
				(1, g_1, \ldots, g_{n+1}), & j = 0 \\
				(\ldots, g_j, 1, g_{j+1}, \ldots), & 0 < j \leq n.
			\end{cases}
		\end{aligned}
	\end{equation*}

	Terdapat funktor alami
	$(\mathcal{E}\Gamma)^{\opp}\to(\mathcal{B}\Gamma)^{\opp}$ yang memetakan
	semua objek ke $\star$ dan setiap morfisme $h\in\Gamma$ ke $h$.
	Morfisme terinduksi $\mathrm{E}\Gamma\to\mathrm{B}\Gamma$ tidak lain
	memetakan $(g_1,\ldots,g_{n+1})$ ke $(g_1,\ldots,g_n)$. Aksi perkalian
	kanan $\Gamma$ pada $(\mathrm{E}\Gamma)_n$ bekerja di sepanjang seratnya, atau,
	dengan kata lain, perubahan data $g_{n+1}$:
	% segment-id: o014.li2.chapter8.nerves-of-categories.q017
	\[ (g_1, \ldots, g_n, g_{n+1})g = (g_1, \ldots, g_n, g_{n+1}g). \]
	Aksi ini jelas berkomutasi dengan $d_i$ dan $s_j$. Jika $\Gamma$ grup,
	aksinya juga bebas.

	Dalam topologi, $\mathrm{B}\Gamma$ (atau realisasi geometrisnya
	$|\mathrm{B}\Gamma|$; lihat \S\ref{sec:geom-realization}) berfungsi
	mengklasifikasikan torsor-$\Gamma$ di atas ruang, sedangkan
	$\mathrm{E}\Gamma\to\mathrm{B}\Gamma$ memberikan torsor universal.
	Contoh~\ref{eg:classifying-contractible} akan menunjukkan bahwa
	$\mathrm{E}\Gamma$ ``dapat dikontraksi dari kanan''.
\end{example}
